% Part: first-order-logic
% Chapter: introduction
% Section: first-order-logic

\documentclass[../../../include/open-logic-section]{subfiles}

\begin{document}

\olfileid{fol}{int}{fol}

\olsection{منطق مرتبهٔ اول}

احتمالاً از نخستین آشنایی خود با منطق صوری، با منطق مرتبهٔ اول
آشنایید.\footnote{در واقع، کمابیش فرض می‌کنیم که چنین است!{} اگر چنین
نیست، می‌توانید کتاب درسی مقدماتی‌تری، مانند \emph{forall x}
\citep{Magnus2021}، را مرور کنید.} شاید آن را با نام «منطق سورها» یا
«منطق محمولات» بشناسید. منطق مرتبهٔ اول، پیش از هر چیز، زبانی صوری
است. یعنی واژگان معینی دارد و عبارت‌هایش رشته‌هایی از همین واژگان‌اند.
اما هر رشته‌ای مجاز نیست. عبارت‌های مجاز انواع گوناگونی دارند: ترم‌ها،
!!{formula}s و !!{sentence}s. توجه اصلی ما به !!{sentence}s منطق مرتبهٔ
اول است: آن‌ها همتایی صوری برای جمله‌های زبان طبیعی به دست می‌دهند و
می‌توانیم پرسش‌هایی را درباره‌شان مطرح کنیم که معمولاً مورد توجه منطق‌دان
است. برای نمونه:
\begin{itemize}
    \item آیا $!B$ منطقاً از~$!A$ نتیجه می‌شود؟
    \item آیا $!A$ منطقاً صادق، منطقاً کاذب، یا ممکن است؟
    \item آیا $!A$ و $!B$ هم‌ارزند؟
\end{itemize}

این پرسش‌ها در وهلهٔ نخست پرسش‌هایی دربارهٔ «معنای» !!{sentence}s منطق
مرتبهٔ اول‌اند. برای نمونه، فیلسوف پرسش از اینکه آیا $!B$ منطقاً
از~$!A$ نتیجه می‌شود را چنین تحلیل می‌کند: آیا حالتی وجود دارد که در آن
$!A$ صادق و~$!B$ کاذب باشد (پس $!B$ از~$!A$ نتیجه نمی‌شود)، یا هر حالتی
که $!A$ را صادق می‌سازد~$!B$ را نیز صادق می‌سازد (پس $!B$ از~$!A$
نتیجه می‌شود)؟ اما هنوز نگفته‌ایم «حالت» چیست---این وظیفهٔ
\emph{معناشناسی} است. معناشناسی منطق مرتبهٔ اول مدلی از تصور شهودی
فیلسوف دربارهٔ «حالت» به دست می‌دهد که از نظر ریاضی دقیق است؛ و نیز---که
مهم است---دقیقاً بیان می‌کند !!a{sentence}~$!A$ چگونه \emph{در یک حالت
صادق} است. مدل دقیق ریاضی‌ای را که خواهیم ساخت !!a{structure} می‌نامیم.
رابطه‌ای را که «صادق در» را دقیق می‌سازد رابطهٔ \emph{ارضا} می‌نامند.
پس آنچه تعریف خواهیم کرد «$!A$ در~$\Struct{M}$ ارضا می‌شود» است (به
نماد: $\Sat{M}{!A}$)، برای !!{sentence}s~$!A$ و
!!{structure}s~$\Struct{M}$. پس از این کار، می‌توانیم اصطلاحات معنایی
دیگر، مانند «از آن نتیجه می‌شود» یا «منطقاً صادق است»، را نیز به‌دقت
تعریف کنیم. این تعریف‌ها امکان می‌دهند با دقت ریاضی تعیین کنیم که، برای
نمونه، آیا
$\lforall[x][(!A(x) \lif !B(x))],
\lexists[x][!A(x)] \Entails \lexists[x][!B(x)]$. پاسخ البته «بله» است.
اگر پیش‌تر نمادگذاری جمله‌های زبان طبیعی در منطق مرتبهٔ اول را آموخته
باشید، این را مثلاً نمادگذاریِ «همهٔ مورچه‌ها حشره‌اند؛ مورچه وجود دارد؛
پس حشره وجود دارد» خواهید شناخت. این استدلال آشکارا معتبر است؛ پس مدل
ریاضی ما از «نتیجه‌شدن» در زبان صوری نیز باید همین پاسخ را بدهد.

موضوع دیگری که احتمالاً از نخستین آشنایی خود با منطق صوری به یاد دارید،
وجود \emph{!!{derivation}s} است. اگر درس مقدماتی منطق صوری گذرانده
باشید، مدرس از شما خواسته است یافتن چنین !!{derivation}sی را تمرین
کنید؛ شاید حتی !!a{derivation} که نشان دهد استلزام بالا برقرار است.
روش‌های گوناگونی برای ارائهٔ !!{derivation}s وجود دارد: ممکن است با
چیزی به نام «استنتاج طبیعی» یا «درخت‌های صدق» کار کرده باشید، اما
روش‌های بسیار دیگری نیز هست. هدف دستگاه‌های !!{derivation} فراهم‌کردن
ابزارهایی است که بتوان با آن‌ها به پرسش‌های بالا پاسخ داد. برای نمونه،
!!a{derivation} استنتاج طبیعی که در آن
$\lforall[x][(!A(x) \lif !B(x))]$ و $\lexists[x][!A(x)]$ مقدمه‌اند و
$\lexists[x][!B(x)]$ نتیجه (سطر پایانی) است، \emph{تأیید می‌کند} که
$\lexists[x][!B(x)]$ منطقاً از $\lforall[x][(!A(x) \lif !B(x))]$ و
$\lexists[x][!A(x)]$ نتیجه می‌شود.

اما چرا؟ در نگاه نخست، دستگاه‌های !!{derivation} هیچ ربطی به معناشناسی
ندارند: ارائهٔ یک !!{derivation} صوری صرفاً مستلزم آن است که نمادها را
به شیوه‌هایی معین و تابع قواعد مرتب کنیم؛ این دستگاه‌ها اصلاً از «حالت»
یا «صادق در» سخن نمی‌گویند. پیوند میان دستگاه‌های !!{derivation} و
معناشناسی باید با پژوهشی فرامنطقی برقرار شود. آنچه لازم است برهانی
ریاضی است؛ برای نمونه، برهانی بر اینکه !!a{derivation} صوری از
$\lexists[x][!B(x)]$ با مقدمات $\lforall[x][(!A(x) \lif !B(x))]$ و
$\lexists[x][!A(x)]$ ممکن است اگر و تنها اگر
$\lforall[x][(!A(x) \lif !B(x))]$ و $\lexists[x][!A(x)]$ با هم
$\lexists[x][!B(x)]$ را نتیجه دهند. اما پیش از آنکه بتوان چنین کرد،
باید کار دقیق و پرزحمتی انجام داد تا تعریف‌های نحو و معناشناسی درست
صورت‌بندی شوند.

\end{document}
